% Options for packages loaded elsewhere
\PassOptionsToPackage{unicode}{hyperref}
\PassOptionsToPackage{hyphens}{url}
\PassOptionsToPackage{dvipsnames,svgnames,x11names}{xcolor}
%
\documentclass[
  11pt,
  a4paper,
]{ltjsarticle}
\usepackage{amsmath,amssymb}
\usepackage{iftex}
\ifPDFTeX
  \usepackage[T1]{fontenc}
  \usepackage[utf8]{inputenc}
  \usepackage{textcomp} % provide euro and other symbols
\else % if luatex or xetex
  \usepackage{unicode-math} % this also loads fontspec
  \defaultfontfeatures{Scale=MatchLowercase}
  \defaultfontfeatures[\rmfamily]{Ligatures=TeX,Scale=1}
\fi
\usepackage{lmodern}
\ifPDFTeX\else
  % xetex/luatex font selection
\fi
% Use upquote if available, for straight quotes in verbatim environments
\IfFileExists{upquote.sty}{\usepackage{upquote}}{}
\IfFileExists{microtype.sty}{% use microtype if available
  \usepackage[]{microtype}
  \UseMicrotypeSet[protrusion]{basicmath} % disable protrusion for tt fonts
}{}
\makeatletter
\@ifundefined{KOMAClassName}{% if non-KOMA class
  \IfFileExists{parskip.sty}{%
    \usepackage{parskip}
  }{% else
    \setlength{\parindent}{0pt}
    \setlength{\parskip}{6pt plus 2pt minus 1pt}}
}{% if KOMA class
  \KOMAoptions{parskip=half}}
\makeatother
\usepackage{xcolor}
\usepackage[margin=24mm]{geometry}
\setlength{\emergencystretch}{3em} % prevent overfull lines
\providecommand{\tightlist}{%
  \setlength{\itemsep}{0pt}\setlength{\parskip}{0pt}}
\setcounter{secnumdepth}{5}
\ifLuaTeX
  \usepackage{selnolig}  % disable illegal ligatures
\fi
\IfFileExists{bookmark.sty}{\usepackage{bookmark}}{\usepackage{hyperref}}
\IfFileExists{xurl.sty}{\usepackage{xurl}}{} % add URL line breaks if available
\urlstyle{same}
\hypersetup{
  colorlinks=true,
  linkcolor={blue},
  filecolor={Maroon},
  citecolor={Blue},
  urlcolor={blue},
  pdfcreator={LaTeX via pandoc}}

\author{}
\date{}

\begin{document}

\hypertarget{pcaux4e3bux6210ux5206ux5206ux6790}{%
\section{PCA（主成分分析）}\label{pcaux4e3bux6210ux5206ux5206ux6790}}

データ行列 \texttt{X∈R\^{}\{n×d\}}
の各行をサンプルとし、まず各特徴量の平均を引いて中心化します。

\[
X_c=X-\mathbf{1}\mu^\top
\]

標本共分散行列は

\[
C=\frac{1}{n-1}X_c^\top X_c
\]

です。

\hypertarget{pca-ux306eux76eeux7684}{%
\subsection{PCA の目的}\label{pca-ux306eux76eeux7684}}

単位ベクトル \texttt{v} 方向へ射影した値は \texttt{X\_cv}
です。その標本分散は

\[
\frac{1}{n-1}\|X_cv\|_2^2
=\frac{1}{n-1}v^\top X_c^\top X_cv
=v^\top Cv
\]

です。

したがって PCA は

\[
\max_{\|v\|=1}v^\top Cv
\]

を解き、最も分散の大きい方向を探します。この解が \texttt{C}
の最大固有値に対応する固有ベクトルです。

\hypertarget{svd-ux3068ux306eux95a2ux4fc2}{%
\subsection{SVD との関係}\label{svd-ux3068ux306eux95a2ux4fc2}}

\[
X_c=U\Sigma V^\top
\]

なら

\[
C=\frac{1}{n-1}V\Sigma^2V^\top
\]

です。

したがって

\begin{itemize}
\tightlist
\item
  \texttt{V} の列 = 主成分方向
\item
  \texttt{σ\_i\^{}2/(n-1)} = 各主成分の標本分散
\end{itemize}

です。

上位 \texttt{k} 主成分への座標は

\[
Z=X_cV_k
\]

です。

\hypertarget{ux6df1ux3044ux898bux65b9}{%
\subsection{深い見方}\label{ux6df1ux3044ux898bux65b9}}

PCA
は「分散最大化」と「再構成誤差最小化」という二つの見方を持ちます。上位主成分を残すことは、中心化データの最良
rank-k 近似を作ることと同値です。

\end{document}
